\documentclass[12pt]{article}
\usepackage[american]{circuitikz}
\usepackage{amsmath}
\usepackage{geometry}
\geometry{margin=2cm}

\begin{document}
\section{Ejemplos de circuitos eléctricos}
\subsection*{Circuito 1: Fuente DC y Resistencia}

\centering
\begin{circuitikz}
\draw
(0,0) to [battery, l=$5v$] (0,3)
    to[R, l=$220 \Omega$] (3,3)
      --(3,0)--(0,0)
\end{circuitikz}


\subsection*{Circuito 2: resistencia y capacitor en serie}

\centering
\begin{circuitikz}
\draw
(0,0) to [battery, l=$5v$] (0,3)
    
    to[R, l=$220 \Omega$] (3,3)
    to[C, l=$4\micro F$] (6,3)
      --(6,0)--(0,0)
\end{circuitikz}





\subsection*{Circuito 3: Inductor y resistencia en serie}

\centering
\begin{circuitikz}
\draw
(0,0) to [battery, l=$5v$] (0,3)
    to[L, l=$5mH $] (3,3)
    to[R, l=$220 \Omega$] (6,3)
      --(6,0)--(0,0)
\end{circuitikz}


\subsection*{Circuito 4: LEDs en paralelo}
\begin{circuitikz}
    \draw
    (0,0) to [battery1, l=$v_s$] (0,4)
    node at (-2,2) {$5v$}
    \draw(0,4)--(4,4)
    \draw(0,0)--(4,0)
    \draw(1,4) to [led, l=LED1] (1,0)
    \draw(4,4) to [led, l=LED2] (4,0)
    
\end{circuitikz}




\end{document}


